\section{Weighted projections and actual range restrictions}

The algebraic identities in this section hold in any unital ring. The
subsequent restrictions are bounded linear maps on the stated normed spaces.
They are supporting identities for the weighted truncation argument
of~\cite[lines 1010--1024]{Arad2016RigorousRG}.

\begin{theorem}[Cancellation on the retained range]
    \label{thm:weighted_right_support_cancellation}
    \lean{mul_inverse_weight_mul_eq_of_right_support}
    \leanok
    If $W_-W=1$ and $D(1-Q)=0$, then $(DW_-)(WQ)=D$.
\end{theorem}
\begin{proof}
    \leanok
    The support identity gives $DQ=D$. Associate the product to cancel
    $W_-W$, leaving $DQ$.
\end{proof}

\begin{theorem}[Factorization of the weighted forcing term]
    \label{thm:weighted_forcing_identity}
    \lean{WeightedProjection.forcing_identity}
    \leanok
    Suppose $WQ=QW$ and $W_-W=1$. For arbitrary $Y,R$,
    \begin{align}
        WQY(1-Q)R=(QWYW_-)(W(1-Q)R).
        \label{eq:weighted_forcing_factorization}
    \end{align}
    The elements $Q,R$ need not be idempotent for this identity.
\end{theorem}
\begin{proof}
    \leanok
    Cancel the adjacent factors $W_-W$ and commute $Q$ past $W$.
\end{proof}

\begin{theorem}[The Sylvester equation between projection ranges]
    \label{thm:weighted_restricted_sylvester}
    \lean{ContinuousLinearMap.exists_restricted_weighted_sylvester}
    \leanok
    Let $E$ be a normed space over a nontrivially normed field, and let
    $X,Y,W,W_-,Q,R$ be bounded endomorphisms of $E$. Suppose $Q^2=Q$,
    $R^2=R$, $W_-W=\Id$, and
    \begin{align}
        WQ & =QW, & XW & =WX, & XQ & =QX, & R(X+Y) & =(X+Y)R.
        \label{eq:weighted_range_commutations}
    \end{align}
    Put $U=\operatorname{ran}Q$, $G=\operatorname{ran}R$,
    $J=QWYW_-Q$ and $L=WQY(1-Q)$. Denote the inclusions by $\iota_U,\iota_G$.
    There are bounded maps $B:G\to G$, $X_U:U\to U$, $Z:G\to U$,
    $V:U\to U$ and $F:G\to U$ satisfying
    \begin{align}
        \iota_G B & =(X+Y)|_G, & \iota_U X_U & =X|_U,\notag                            \\
        \iota_U Z & =(WQ)|_G,  & \iota_U V   & =J|_U,       & \iota_U F & =L|_G,\notag \\
        X_UZ-ZB   & =-VZ-F.
        \label{eq:weighted_range_sylvester}
    \end{align}
    Orthogonality of $Q,R$ is not required.
\end{theorem}
\begin{proof}
    \leanok
    The commutations in~(\ref{eq:weighted_range_commutations}) imply that
    $X+Y$ preserves $G$ and $X$ preserves $U$. The ranges of $WQ$, $J$
    and $L$ lie in $U$. Restrict these actual maps. Before restriction,
    associativity, idempotence and the commutations give
    \begin{align}
        X(WQR)-(WQR)(X+Y)=-J(WQR)-LR.
        \label{eq:weighted_ambient_sylvester}
    \end{align}
    Since $R$ fixes $G$, restricting~(\ref{eq:weighted_ambient_sylvester})
    yields~(\ref{eq:weighted_range_sylvester}).
\end{proof}

\begin{theorem}[A lower quadratic bound on the projection range]
    \label{thm:projection_range_quadratic_lower}
    \lean{ContinuousLinearMap.exists_restrict_range_quadratic_lower}
    \leanok
    Let $E$ be a complex inner-product space, and let $Q,X$ be bounded
    endomorphisms with $Q^2=Q$, $XQ=QX$ and $uQ\le X$, where $u\in\R$.
    There is a bounded map $X_U:U\to U$, with $U=\operatorname{ran}Q$,
    such that $\iota_U X_U=X|_U$ and
    $u\|y\|^2\le\operatorname{Re}\langle y,X_Uy\rangle$ for every $y\in U$.
    The idempotent $Q$ need not be orthogonal.
\end{theorem}
\begin{proof}
    \leanok
    Commutation makes $U$ invariant. For $y\in U$, idempotence gives
    $Qy=y$, so positivity of $X-uQ$ gives the asserted bound.
\end{proof}

\begin{theorem}[Norm from the ambient realization]
    \label{thm:restricted_map_ambient_norm}
    \lean{ContinuousLinearMap.norm_le_of_subtypeL_comp_eq_domRestrict}
    \leanok
    Let $G,U$ be subspaces of a normed space $E$ over a nontrivially
    normed field. If bounded maps $f:G\to U$ and $A:E\to E$ satisfy
    $\iota_Uf=A|_G$, then $\|f\|\le\|A\|$.
\end{theorem}
\begin{proof}
    \leanok
    The inclusion $\iota_U$ is isometric, and restriction to $G$ cannot
    increase the operator norm. This also covers zero-dimensional subspaces.
\end{proof}

\begin{theorem}[Order under the Euclidean matrix action]
    \label{thm:matrix_euclidean_order}
    \lean{Matrix.toEuclideanCLM_le_of_le}
    \leanok
    If complex square matrices satisfy $A\le B$, their bounded actions on
    the corresponding finite-dimensional Euclidean space satisfy the same order.
    The index set may be empty.
\end{theorem}
\begin{proof}
    \leanok
    The positive matrix $B-A$ defines a positive Euclidean operator.
\end{proof}
